% Aqui começa o capítulo de Introdução.
% Use o comando \label para definir um rótulo, 
% caso seja necessário referenciar esse capítulo
% posteriormente.
\chapter{Introdução}\label{chp:Introducao}

O câncer de próstata configura-se como uma das neoplasias malignas de maior prevalência entre indivíduos do sexo masculino, sendo o quarto em incidência global e o oitavo em letalidade \cite{Ferlay_GCO_2024}. A identificação precoce da patologia é fundamental para o êxito no tratamento, contribuindo para a diminuição da mortalidade e para a melhora nos desfechos clínicos dos pacientes. Entretanto, os métodos diagnósticos atualmente empregados apresentam limitações, incluindo suscetibilidade a erros humanos, o que pode resultar em casos em que a doença não é detectada corretamente ou em diagnósticos incorretos de câncer em pacientes saudáveis, ambos com implicações clínicas relevantes \cite{Beltran2019}.

A aplicação de modelos de aprendizagem profunda (do inglês \Sigla{\textit{Deep Learning}}{DL}) no diagnóstico de câncer de próstata oferece uma oportunidade promissora para melhorar a acurácia e a consistência dos diagnósticos. Esses modelos podem auxiliar patologistas na análise de imagens histológicas, reduzindo o impacto de fatores como fadiga e variações na prática da observação microscópica \cite{Perincheri2021}.

 Com o advento da DL, a digitalização de WSIs se consolidou como realidade clínica e de pesquisa. WSIs são geradas por scanners de alta resolução que criam arquivos digitais da lâmina inteira, potencialmente gerando imagens na ordem de gigabytes por exame \cite{Pantanowitz2011}. Esse volume massivo de dados impõe graves desafios computacionais relacionados ao armazenamento, à transmissão e ao processamento dessas imagens, já que os padrões teciduais são complexos e de alta heterogeneidade \cite{Gurcan2009}.

Além disso, a análise manual dessas imagens por patologistas, embora ainda essencial, revela-se extremamente laboriosa e demorada na prática. O fluxo de trabalho tradicional requer percorrer manualmente Regiões de Interesse (do inglês \Sigla{\textit{Regions of Interest}}{ROI}) em lâminas inteiras, o que consome tempo e recursos humanos significativos \cite{SchomigMarkiefka2021}. A essa sobrecarga soma-se a variabilidade interobservador: diferentes patologistas podem chegar a conclusões distintas ao interpretar as mesmas regiões teciduais, o que impacta diretamente na reprodutibilidade e na acurácia diagnóstica \cite{CAP2018}.

A \Sigla{Inteligência Artificial}{IA}, e mais especificamente DL, emergiu como uma ferramenta promissora para auxiliar na análise de WSIs, permitindo automatizar tarefas antes restritas à interpretação manual de patologistas \cite{Litjens2017}. Modelos de DL, particularmente CNNs, têm demonstrado capacidade notável na identificação de padrões complexos em imagens médicas, como a detecção e segmentação de regiões cancerosas em lâminas histológicas \cite{Ciresan2012}. Mais recentemente, arquiteturas baseadas em \textit{Transformers}, originalmente desenvolvidas para \Sigla{Processamento de Linguagem Natural}{PLN}, vêm sendo adaptadas para visão computacional, capturando relacionamentos globais em WSIs que complementam as características locais extraídas pelas CNN \cite{Dosovitskiy2021}.

Apesar do grande volume de pesquisas focadas no desenvolvimento e aprimoramento de modelos de DL para aplicações envolvendo câncer de próstata \cite{Rabilloud} , observa-se uma lacuna na literatura no que diz respeito à criação e validação de metodologias ou \textit{framework} computacionais completos e replicáveis. A falta de um processo metodológico padronizado e bem documentado dificulta a comparação entre diferentes abordagens e a adoção mais ampla dessas tecnologias na prática clínica ou pesquisas futuras.

Esta dissertação aborda essa lacuna ao propor o desenvolvimento e a validação de um \textit{framework} computacional \textit{end-to-end} para a segmentação semântica de tecido canceroso em WSIs de câncer de próstata. O objetivo central não é apenas alcançar alta performance na tarefa de segmentação, mas também estabelecer uma metodologia sistemática e replicável que possa servir de base para futuros trabalhos na área, sendo independente da arquitetura específica utilizada dentro do \textit{framework}. Parte desta pesquisa foi desenvolvida em colaboração com uma equipe de pesquisadores da Universidad de La Frontera, Chile\footnote{Universidad de La Frontera: \url{https://www.ufro.cl/}}, enriquecendo a abordagem e os dados utilizados.

Este capítulo introdutório contextualiza o problema, destaca a relevância clínica e computacional da detecção de câncer de próstata em patologia digital, aponta as limitações das abordagens atuais e apresenta a proposta central desta dissertação. As seções subsequentes detalharão os objetivos, a justificativa e a estrutura completa do documento.

\section{Objetivos}\label{sec:objetivos}

O objetivo geral desta dissertação é desenvolver e validar um \textit{framework} computacional, \textit{end-to-end}, que integre CNNs e \textit{Transformers} para a segmentação semântica precisa de tecido canceroso em WSIs de biópsias de câncer de próstata, estabelecendo assim uma metodologia modular, documentada e replicável para pesquisas futuras na área.

Para alcançar o objetivo geral, os seguintes objetivos específicos são buscados:
\begin{itemize}
    \item \textbf{Desenvolver e Validar Módulos de Pré-processamento de WSIs:} Projetar e implementar os módulos iniciais do \textit{framework} para a manipulação eficiente de WSIs, incluindo técnicas para leitura, extração de \textit{patches} e segmentação de tecido. A validação deste objetivo focará na robustez dos métodos em diferentes cenários de imagem.

    \item \textbf{Analisar o Impacto da Curadoria e Padronização de Dados:} Quantificar sistematicamente como a qualidade dos dados de treinamento influencia o desempenho do \textit{framework}, investigando o impacto da remoção de artefatos, da filtragem de \textit{patches} com anotações ambíguas ou incorretas (como regiões de fundo anotadas como tecido), e comparando o efeito de diferentes métodos de normalização de cor na robustez e acurácia final da segmentação.

    \item \textbf{Investigar a Eficácia de Ensembles Híbridos para Segmentação:} Construir, treinar e otimizar \textit{ensembles} de modelos compostos por uma diversidade de arquiteturas, incluindo tanto CNNs quanto \textit{Transformers}. O objetivo é explorar a sinergia entre a capacidade das CNNs de capturar características locais e a habilidade dos \textit{Transformers} em modelar dependências globais, visando obter um resultado de segmentação mais robusto e preciso do que qualquer modelo individualmente. O desempenho final do \textit{ensemble} será avaliado de forma quantitativa e qualitativa em um conjunto de teste independente.
    
    \item \textbf{Validar a Generalização e Transferibilidade do \textit{Framework}:} Avaliar a aplicabilidade e a robustez da metodologia \textit{end-to-end} proposta em cenários que vão além do conjunto de dados principal. Isso inclui testar o \textit{framework} em outros contextos histopatológicos (como outros tipos de câncer), investigando sua capacidade de manter um desempenho consistente e demonstrando a sua generalidade como uma solução metodológica.
\end{itemize}

\section{Justificativa}\label{sec:justificativa}
A relevância desta pesquisa assenta-se em múltiplos pilares, abrangendo aspectos clínicos, computacionais e científicos.

A importância clínica do câncer de próstata é inegável. Sendo uma das maiores causas de mortalidade por câncer em homens, a melhoria nos métodos diagnósticos tem impacto direto na saúde pública. Ferramentas de diagnóstico auxiliadas por computador (do inglês \Sigla{\textit{Computer-Aided Diagnosis}}{CAD}) baseadas em DL têm o potencial de otimizar o fluxo de trabalho dos patologistas, reduzindo tempo de análise, minimizando a fadiga visual e potencialmente aumentando a consistência e acurácia diagnóstica, especialmente em casos limítrofes ou complexos. A segmentação precisa do tecido canceroso é um passo fundamental para quantificar a extensão do tumor e auxiliar na escala de Gleason \cite{Gleason1966}, que é informação essencial para o planejamento terapêutico do paciente.

Do ponto de vista computacional, trabalhar com WSIs representa um desafio significativo. O tamanho das imagens exige estratégias eficientes de processamento e análise que demandam considerável poder computacional e expertise técnica. Desenvolver uma metodologia que seja replicável, utilizando recursos como \textit{Google Colab} e bibliotecas otimizadas como \textit{PyTorch}, torna a pesquisa em patologia digital mais acessível e viável para um número maior de grupos de pesquisa, democratizando o avanço na área.

Cientificamente, existe uma lacuna metodológica identificada. Há certa escassez de estudos que se dediquem a estabelecer e documentar um \textit{framework} completo e replicável . A ciência progride através da replicabilidade e da construção sobre trabalhos anteriores. Dessa forma, um \textit{framework} bem definido, que detalhe desde o pré-processamento dos dados até a avaliação final, e que seja agnóstico a estilos específicos de arquiteturas de DL, serve como um alicerce robusto. Ele permite que outros pesquisadores reproduzam os resultados, testem novas hipóteses dentro de um ambiente controlado e comparem resultados de forma justa. A proposta de integrar CNNs e \textit{Transformers} dentro deste \textit{framework} também se alinha às tendências atuais da pesquisa em visão computacional, explorando sinergias entre diferentes paradigmas de modelagem para tarefas complexas de imagem.

Nesse contexto, a criação de uma estrutura \textit{end-to-end} replicável, conforme proposta nesta dissertação, não apenas oferece uma ferramenta potencialmente útil para a análise de câncer de próstata, mas também estimula novos avanços na área ao fornecer uma base metodológica sólida e transparente, fomentando a colaboração e acelerando o ciclo de inovação em patologia digital computacional.

\section{Estrutura do documento}\label{sec:estrutura_documento}
Esta dissertação está organizada da seguinte forma, visando apresentar de maneira clara e lógica o desenvolvimento e os resultados da pesquisa:

\textbf{Capítulo 1: Introdução:} Apresenta o contexto geral do câncer de próstata e da patologia digital, introduz os desafios relacionados às WSIs e ao diagnóstico, destaca o potencial do DL (CNNs e \textit{Transformers}), cita a lacuna metodológica na literatura e delineia a proposta central da dissertação. Inclui também os objetivos, a justificativa e esta descrição da estrutura do documento.

\textbf{Capítulo 2: Levantamento bibliográfico:} Revisa os conceitos fundamentais necessários para a compreensão do trabalho. Abrange noções sobre câncer de próstata, patologia digital e WSIs, princípios de DL, arquiteturas de CNNs e \textit{Transformers}, e técnicas de segmentação semântica. Apresenta também uma revisão dos trabalhos relacionados na área de detecção e segmentação de câncer de próstata em WSIs usando IA, comparando metodologias e resultados existentes.

\textbf{Capítulo 3: Desenvolvimento:} Apresenta, de forma detalhada, o \textit{framework} \textit{end-to-end} desenvolvido nesta pesquisa. São descritas as fontes de dados empregadas, incluindo os processos de aquisição, organização e anonimização, bem como as estratégias adotadas para o pré-processamento das WSIs — abrangendo a geração de \textit{patches}, os métodos de identificação e exclusão de artefatos, e os procedimentos de correção de imprecisões nas anotações patológicas dos casos de câncer. Além disso, discute-se as técnicas de \textit{data augmentation} utilizadas, as arquiteturas de modelos baseadas em CNNs e \textit{Transformers} avaliadas, o pipeline de treinamento (funções de perda, otimizadores, particionamento em treino/validação/teste) e as métricas selecionadas para avaliação do desempenho. O capítulo também descreve as principais ferramentas e bibliotecas computacionais utilizadas, tais como \textit{Python}, \textit{PyTorch} e \textit{Google Colab Pro}.


\textbf{Capítulo 4: Resultados e discussão:} Apresenta os experimentos realizados para validar a estrutura e avaliar o desempenho do \textit{framework} em bases de dados diversas. Descreve a configuração experimental, exibe os resultados quantitativos obtidos (tabelas e gráficos com as métricas de avaliação) e fornece exemplos qualitativos das segmentações geradas pelos modelos nas WSIs. Apresenta uma análise crítica dos resultados, comparando diferentes abordagens testadas dentro do \textit{framework} e discutindo os achados à luz dos objetivos propostos e da literatura.